%! Author = taj
%! Date = 7/22/25

% Preamble
\documentclass[11pt]{article}

% Packages
\usepackage{amsmath}
\usepackage{graphicx}
\usepackage{float}

\title{1. Merjenje hitrosti tekočin z Venturjevo cevjo}
\author{Taj Šobot}

% Document
\begin{document}
    \maketitle
    V tem eksperimentu smo s pomočjo Venturjeve cevi merili hitrost, masni pretok in tlačno razliko v ožjem in širšem delu cevi.


    \section{Meritve}
    Izmerjeni zunanji premer ($err_d/D = 0.01cm$):
    \[
        D_1 = 3.13 cm; D_2 = 2.98 cm; D_3 = 3.16 cm; D_4 = 2.95 cm.
    \]
    Izmerjeni notranji premeri
    \[
        d_1 = 2.69 cm; d_2 = 2.73 cm; d_3 = 2.77 cm; d_4 = 2.83 cm.
    \]
    Višinska razlika teločine v u cevki (pri dveh različnih nastavitvah hitrosti):
    $err_h = 0.1cm$
    \begin{table}[h!]
        \centering
        \begin{tabular}{|c|c|}
            \hline
            \textbf{$\Delta h_{a} [cm]$} & \textbf{$\Delta h_b$ [cm]} \\
            \hline
            3.0  & 4.9 \\
            3.1  & 4.8 \\
            3.1  & 4.8 \\
            3.1  & 4.8 \\
            3.1  & 4.8 \\
            3.1  & 4.7 \\
            \hline
        \end{tabular}
        \caption{tabela meritev u cevke}
        \label{tab:dh_data}
    \end{table}

    Skica eksperimenta:
    \begin{figure}[H]
        \centering
        \includegraphics[width=0.5\textwidth]{img}
        \caption{Venturjeva cev z oznakami}
        \label{fig:2}
    \end{figure}

    \section{rezultati in grafi}
    Iz premerov lahko določimo povprečni premer:
    \begin{align*}
        D &= (3.1 \pm 0.1)\,\text{cm} \\
        d &= (2.8 \pm 0.1)\,\text{cm}
    \end{align*}
    Iz tabele določimo razliko višine v u cevki:
    \begin{align*}
        h_a &= (3.1 \pm 0.1)\,\text{cm} \\
        h_b &= (4.8 \pm 0.1)\,\text{cm}
    \end{align*}

    \subsection{hitrost}
    In ploščine $S$ preko enačbe $S_1 =\frac{1}{4} \pi D^2 $ in $S_2 =\frac{1}{4} \pi d^2 $ in dobimo:
    \begin{align*}
        S_1 = (7.5 \pm 0.3) cm^2 \\
        S_2 = (6.2 \pm 0.3) cm^2
    \end{align*}

    nato vstavimo v enačbo($\rho_v = 1.00 \frac{g}{cm^3}$, $\rho_z = 1.23*10^{-3} \frac{g}{cm^3}$ in $g = 981 \frac{cm}{s^2}$):
    \[
        v_{1a} = \sqrt{\frac{2 \rho_v g h_a }{\rho_z[(\frac{S_1}{S_2})^2 -1]}} = (33 \pm 4) \frac{m}{s}
    \]
    Hitrost v zoženem delu pa dobimo preko enačbe $v_2 = v_1 \frac{S_1}{S_2}$:
    \[
        v_{2a} = (40 \pm 4) \frac{m}{s}
    \]

    In še za b) primer:
    \[
        v_{1b} = \sqrt{\frac{2 \rho_v g h_b }{\rho_z[(\frac{S_1}{S_2})^2 -1]}} = (41 \pm 4) \frac{m}{s}
    \]
    Ter hitrost v zoženem delu:
    \[
        v_{2b} = (50 \pm 5) \frac{m}{s}
    \]
    \subsection{tlačna razlika}
    Tlačna razlika se izračuna z enačbo:
    \[
        \Delta p = \rho_v g h
    \]
    Toerj:
    \begin{align*}
        \Delta p_a = (300 \pm 20) Pa \\
        \Delta p_b = ( 470 \pm 30) Pa
    \end{align*}

    \subsection{masni pretok}
    Masni pretok se izračuna s sledečo enačbo(za oba dela cevi more bit isto):
    \[
        \Phi_{1} = \Phi_{2} = S_1 v_1 = S_2 v_2
    \]
    Torej a):
    \begin{align*}
        \Phi_{1a} = (25 \pm 2 ) \frac{g}{s} \\
        \Phi_{2a} = (25 \pm 2) \frac{g}{s}
    \end{align*}

    In b)
    \begin{align*}
        \Phi_{1b} = (30 \pm 2 )\frac{g}{s}  \\
        \Phi_{2b} = (31 \pm 2) \frac{g}{s}
    \end{align*}

    Rezultata se ujemata.
\end{document}